\section{Data}

\subsection{Source}

The empirical analysis uses microdata from the Gran Encuesta Integrada de
Hogares (GEIH), Colombia's nationally representative labor force survey produced
by DANE. The current harmonized file covers the years 2008--2019 and 2021--2025.
The working sample contains 5,115,921 worker-year observations with valid real
hourly labor income and positive expansion weights. The baseline regression
sample contains 5,034,706 observations after requiring non-missing values for the
controls used in the econometric specification.

\subsection{Sample Construction}

The sample is restricted to employed workers with positive labor income, valid
weekly hours, and positive survey expansion factors. Monthly labor income is
taken from the GEIH labor income variable (\texttt{INGLABO}). We convert it to
annual labor income and divide by annual hours, computed as usual weekly hours
times 52:

\[
  w_{it}^{nominal} =
  \frac{12 \times INGLABO_{it}}{52 \times hours_{it}}.
\]

We then express hourly labor income in constant 2025 pesos using the December CPI
series encoded in the data construction scripts:

\[
  w_{it}^{real} = w_{it}^{nominal} \times
  \frac{CPI_{2025}}{CPI_t}.
\]

The main outcome is the logarithm of real hourly labor income,
\(\ln(w_{it}^{real})\). Observations with non-positive income, invalid hours, or
non-positive weights are excluded.

\subsection{Firm Size and Controls}

Firm size is measured using the worker-reported firm-size variable in GEIH. The
construction harmonizes the older \texttt{P6870} categories and the newer
\texttt{P3069} categories into the following ordered groups: Solo, 2-3, 4-5,
6-10, 11-19, 20-30, 31-50, 51-100, and 101+. When the newer questionnaire
separates 101-200 and 201+ workers, both categories are collapsed into 101+ to
maintain comparability with earlier years.

The control variables are also harmonized across survey changes. Education is
grouped into comparable categories: no education, preschool, primary, lower
secondary, upper secondary, and higher education. Formality is defined using
pension contribution status: workers who contribute to pensions are classified
as formal, while workers who do not contribute are classified as informal.
Economic activity is harmonized across CIIU revisions into broad sectors, which
allows the regressions to absorb sector-year shocks and changes in the sectoral
composition of employment.

Two caveats are important. First, GEIH is a household survey, not linked
employer-employee administrative data. Firm size is therefore reported by
workers, and the data do not identify firms. Second, the analysis documents
conditional wage differences by firm size; it does not identify the causal effect
of moving a worker from a small firm to a large firm.
